\chapter{文章を数にする}\label{ch:text}

決算発表文は文章である。回帰や BLM に入れるには、文章を数（点数やベクトル）に変える必要がある。この章では、この研究で比べた 6 つの方法を、古いものから順に説明する。辞書・TF-IDF・FinBERT をまとめて「従来の NLP」、gpt-4o-mini と ChronoGPT を「LLM」と呼ぶ。

\section{トークン}

コンピュータは文章を\term{トークン}の列として扱う。古い方法では単語がトークンで、新しいモデルでは単語をさらに細かく分けた部分（サブワード）がトークンである。英語では、1 トークンはおよそ 4 文字、0.75 語に当たる。ChronoGPT の文脈の上限 1,792 トークンは、英語の約 7,000 字に当たる。

\section{辞書の方法}

あらかじめ「肯定の語」「否定の語」のリストを作っておき、文章の中の語を数える。

\begin{equation}
\text{net\_tone} = \frac{\#\text{肯定} - \#\text{否定}}{\#\text{肯定} + \#\text{否定}}
\end{equation}

一般の感情の辞書は、金融の文章では誤る。「liability（負債）」「tax（税）」「cost（費用）」は一般には否定語だが、決算の文章では中立の会計の用語である。Loughran と McDonald（2011）は、10-K の文章から金融用の辞書を作った。肯定約 350 語、否定約 2,300 語、不確かさ約 300 語などからなる。

\begin{description}[leftmargin=1.5em,style=nextline]
\item[強み] 透明で、何が点数を動かしたかがわかる。計算が速い。学習が要らない。
\item[弱み] 語の順序と否定（「not good」）を見ない。「売上の減少が予想より小さかった」のような文脈を読めない。
\end{description}

\section{TF-IDF}

\term{TF-IDF} は、文章を「語の重みのベクトル」にする。語 $w$ の文章 $d$ での重みは
\begin{equation}
\text{tfidf}(w,d) = \text{tf}(w,d)\times\log\frac{N}{\text{df}(w)}
\end{equation}
である。$\text{tf}$ はその文章での出現の頻度、$\text{df}$ はその語を含む文章の数、$N$ は文章の総数である。どの文章にも出る語（the、company）は重みが小さく、特定の文章にだけ出る語は重みが大きい。

TF-IDF のベクトルそのものは点数ではない。語彙が 2 万なら 2 万次元のベクトルである。これを入力にして、\textbf{過去の発表の反応でリッジ回帰を学習}すると、「どの語が出ると反応が大きいか」をデータから学んだ予測になる。辞書は人が決めた語のリストだが、TF-IDF＋リッジはデータが語の重みを決める。

\section{BERT と FinBERT}

\subsection{Transformer}

2017 年の\term{Transformer} は、文章の各トークンが、ほかのすべてのトークンを参照して自分の表現を作り直す仕組み（\term{注意機構}、アテンション）を何層も重ねる。「not」が「good」の意味を反転させる、「decline」が何の減少かを前後から読む、といった文脈の処理ができる。

\subsection{BERT}

\term{BERT}（2018）は、Transformer に大量の文章を読ませ、穴埋め（隠した語を前後から当てる）で\term{事前学習}したモデルである。事前学習で言語の一般的な知識を身につけ、その後、少ないデータで個別の課題に\term{追加学習}（ファインチューニング）する。入力の長さは 512 トークンまでである。

\subsection{FinBERT}

\term{FinBERT}（\code{yiyanghkust/finbert-tone}）は、BERT を金融の文章（10-K、決算説明会、アナリストのレポート）で追加の事前学習し、さらにアナリストのレポートの 1 万文に付けた「肯定・中立・否定」で学習した分類のモデルである。1 文を入れると、3 つのラベルの確率を返す。

\begin{iobox}[入力と出力：FinBERT の使い方 2 通り]
\iorow{学習なし}{発表文を文に分け、先頭 200 文を 1 文ずつ入れて、肯定の確率 − 否定の確率を平均する（実験 2・11）}
\iorow{追加学習}{先頭 512 トークンを入れ、分類の頭を回帰の頭（数 1 つ）に替えて、\textbf{反応を正解にして全体を学習}する（実験 4・11）}
\end{iobox}

\section{生成型の LLM}

\subsection{次のトークンの予測}

GPT 系の\term{LLM}（大規模言語モデル）は、Transformer の「前だけを見る」形（デコーダ）を、ウェブの大量の文章で「次のトークンを当てる」ように事前学習したものである。パラメータの数は数十億〜数兆で、さらに人の指示に従うように追加学習（指示の学習、人の評価による強化学習）されている。

LLM に課題をさせるには、指示と入力を\term{プロンプト}として渡し、続きを生成させる。例を見せずに指示だけで答えさせることを\term{ゼロショット}と呼ぶ。この研究の gpt-4o-mini はすべてゼロショットで、学習はさせていない。

\subsection{答えの設定}

\begin{description}[leftmargin=1.5em,style=nextline]
\item[温度（temperature）] 次のトークンを選ぶときのランダムさ。0 なら最も確率の高いものを選ぶ（ほぼ決定的）。1 なら確率どおりに選ぶ。実行 1 は 1.0 で 30 個の回答を集め、そのばらつきを $\Omega$ に使った。H06・実験 2 は 0。
\item[構造化出力] 答えを決まった形の JSON に限る機能。各項目の型と値の範囲（$-2$〜$+2$ の整数など）を指定でき、答えの形の誤りがなくなる。
\item[確率（logprobs）] 各トークンの選ばれる確率を返す機能。「1〜5 のどれか」を聞いたとき、5 つの数字の確率から期待値を計算できる（実験 G・H05）。
\item[複数の回答（n）] 1 回の入力で $n$ 個の回答を独立に生成する。入力の費用は 1 回分で済む（実行 1 の $n=30$）。
\end{description}

\subsection{知識の締め切り}

LLM は、事前学習した時点までの文章しか知らない。gpt-4o-mini の知識は 2023 年 10 月までである。それより前の決算発表を読ませると、LLM はその後に起きたこと（株価の動きやニュース）を知っている可能性がある。これを\term{先読み}（ルックアヘッド）の問題と呼ぶ（\chapref{ch:pit}）。

\section{ChronoGPT：時点の決まった LLM}

\term{ChronoGPT}（He・Lv・Manela・Wu、2025。この研究で使ったのは指示に従うよう学習した ChronoGPT-Instruct）は、先読みの問題を解くために作られた LLM である。各年末について、\textbf{その年の 12 月 31 日までに書かれた文章だけ}で学習した\term{版}が公開されている。1.55B パラメータ、文脈 1,792 トークン、英語のみ、MIT ライセンス。この研究では 2014〜2024 年末の 11 個の版を取得した。版 2015 は、2015 年末までの文章しか知らない。2016 年の決算を版 2015 で読めば、先読みは起こりえない。判断日 $d$ の文章は、$d$ の前の年末の版で読む。

この研究では、ChronoGPT を 3 通りに使った。

\begin{longtable}{P{3.0cm}P{6.0cm}P{5.4cm}}
\toprule
\textbf{使い方} & \textbf{中身} & \textbf{結果} \\
\midrule
\endhead
\bottomrule
\endfoot
選択肢の確率（実験 2） & 「ガイダンスは引き上げか」などを選択肢で聞き、各選択肢の確率を較正して点数にする。学習なし & 弱い（実験 5 で埋め込みと同時に +6bp、$t$ 1.2） \\
埋め込み＋リッジ（実験 5・9・11） & 最後の層の出力をトークンで平均した 1,536 個の数を、リッジで反応に回帰する。LLM は学習しない & +34bp。従来の NLP を超える \\
LoRA（実験 10） & LLM の内部に少数の学習できる重みを足し、反応で学習する & +29.5bp。リッジより弱い \\
\end{longtable}

\section{埋め込み}

LLM の各層は、各トークンについて数千個の数のベクトル（\term{隠れ状態}）を出す。これを文章全体で 1 つのベクトルにまとめたものを\term{埋め込み}と呼ぶ。まとめ方にはいくつかある。

\begin{itemize}
\item \textbf{全トークンの平均}（mean pooling）：最後の層の隠れ状態を平均する。実験 5 の方法。
\item \textbf{最後のトークン}：デコーダ型の LLM では、最後のトークンが前のすべてを見ている。指示文（「この発表の要点は：」など）を後ろに付けて、その最後のトークンを使う。
\item \textbf{中間の層}：最後の層は「次のトークンを当てる」ことに特化している。中間の層のほうが、文章の意味を一般的な形で持っているという報告がある。実験 11 の探索では、52 層のうち 26 層目の平均が、最後の層の平均より強かった（$R^2$ 増分 0.49 → 0.80）。
\end{itemize}

埋め込みはそれ自体では点数ではない。「この 1,536 次元のどの方向が反応と関係するか」は、リッジなどで学習する必要がある。LLM が言語の知識を持ち、リッジがその知識のうち株価に関係する部分を拾う、という分業である。

\section{長い文章をどう読むか}

決算発表文は長い。H06 で gpt-4o-mini に渡した入力は 1 本あたり平均約 3,950 トークンで、先頭 15,000 字の上限の近くまで使っていた。モデルには読める長さの上限がある。

\begin{longtable}{P{3.4cm}P{3.6cm}P{7.4cm}}
\toprule
\textbf{モデル} & \textbf{読む範囲} & \textbf{読み残し} \\
\midrule
\endhead
\bottomrule
\endfoot
辞書・TF-IDF & 全文 & なし \\
FinBERT（学習なし） & 先頭 200 文 & 長い発表の後半（多くは表と定型文） \\
FinBERT（追加学習） & 先頭 512 トークン（約 2,000 字） & 本文の大部分 \\
ChronoGPT & 先頭 1,792 トークン（約 7,000 字） & 後半。実験 11 の確認では、修正 1 の後でも、見通しの語が読む範囲の外にある発表が 10.7\% \\
gpt-4o-mini & 先頭 15,000 字 & 長い発表の末尾 \\
\end{longtable}

読み残しを減らす方法として、実験 11 で次を試した。
\begin{itemize}
\item \textbf{掃除}（修正 1）：表の切れ端、免責事項、説明会の案内、会社の紹介を除く。平均 27\% の文字が減った。
\item \textbf{先頭と末尾}（修正 2）：先頭だけでなく末尾も読む。
\item \textbf{圧縮}（修正 5）：\term{LLMLingua-2} は、各トークンを「残す・捨てる」に分類するモデル（XLM-RoBERTa）で、文章を指定の長さまで縮める。約 1,700 トークンに圧縮して ChronoGPT に読ませた。
\end{itemize}
どれも、説明力をほとんど上げなかった（圧縮はむしろ悪くした）。

\section{方法の比較}

\begin{longtable}{P{2.6cm}P{2.8cm}P{2.6cm}P{2.4cm}P{3.6cm}}
\caption{文章を数にする方法の比較（決算発表への 2 日の反応に対する、統制に加えたときの $R^2$ の増分。標本と期間が違うので、同じ行の中でだけ比べる）。}\label{tab:text-compare}\\
\toprule
\textbf{方法} & \textbf{反応で学習} & \textbf{先読み} & \textbf{実験 4・5・10（標本 A）} & \textbf{実験 11（最終の期間 B）} \\
\midrule
\endfirsthead
\toprule
\textbf{方法} & \textbf{反応で学習} & \textbf{先読み} & \textbf{実験 4・5・10（標本 A）} & \textbf{実験 11（最終の期間 B）} \\
\midrule
\endhead
\bottomrule
\endfoot
辞書 & なし & なし & — & +0.02 \\
FinBERT（学習なし） & なし & 2019 年までの文章 & — & +0.24 \\
TF-IDF＋リッジ & リッジ & なし & — & +0.37 \\
ChronoGPT＋リッジ & リッジ & なし（版 2015／2022） & +0.31 & +0.59 \\
ChronoGPT＋LoRA & LoRA & なし（版 2022） & +0.13 & — \\
FinBERT（追加学習） & 全体 & 2019 年までの文章 & +0.84 & +1.01 \\
gpt-4o-mini & なし & 2023-10 まで & +2.01 & +2.49（2023-11 以降） \\
\end{longtable}

学習なしの gpt-4o-mini が、反応で学習したどのモデルよりも強い。ただし、gpt-4o-mini は最も大きく、最も新しい文章まで学習している。強さの理由が「大きさ」なのか「知識の新しさ」なのかは、この研究では分けられない。

\begin{projbox}
\begin{itemize}
\item \textbf{gpt-4o-mini の質問}（H06 のまま、すべての実験で同じ）：発表文の先頭 15,000 字を渡し、\code{guidance\_change}（ガイダンスの変化、$-2$〜$+2$、記載なしは null）、\code{outlook}（今後の事業の見通し、$-2$〜$+2$）、\code{tone}（経営陣の語調、$-2$〜$+2$）、\code{risk\_flags}（10 種類のリスクの印）を構造化出力で答えさせる。合成点数 ＝ tone ＋ outlook ＋ guidance\_change（null は 0）。学習しない。
\item \textbf{点数の付け方の番号}：DB の \code{release\_score\_runs} に、方法ごとのプロンプト・設定・版を JSON で記録した（1 gpt-4o-mini、2 辞書、3 FinBERT、6 ChronoGPT の選択肢、7 TF-IDF、8 FinBERT の追加学習、9 ChronoGPT の LoRA）。
\item \textbf{質問の中立性}：システムプロンプトで「発表文の中身だけで判断し、会社について知っていることや、その後の株価を使わない。どの発表にもある宣伝の言い回しではなく中身を見る」と指示した。
\end{itemize}
\end{projbox}
